\begin{abstract}
\noindent\textbf{Questions.} Which food crops are most vulnerable to observed and projected heat and water stress,
once we count how much people depend on them for nutrients? And if such a crop's yield falls, what happens to a
country's nutrient supply?

\noindent\textbf{Method.} The Climate--Nutrition Vulnerability Index pairs each stress with its own crop response:
$\Fd = \tfrac12(E_{\text{heat}}S_T + E_{\text{water}}S_P)\,\tilde I$, from the observed 2001--2024 stress level
($E$), literature sensitivity (yield lost per degree above a threshold, and FAO's $K_y$; $S$) and the crop's mean
share of national energy, protein, fat, zinc and iron supply ($\tilde I$). It covers nine crops in \nPairs\
crop--country pairs and three periods: 2001--2024, 1971--2000, and 2041--2060 under one climate model (MRI-ESM2-0,
SSP2-4.5) with a fixed normalisation. Nutrient consequences are computed for a standard 10\% yield loss and, as a
\emph{conditional screening estimate}, for the projected climate-driven yield change.

\noindent\textbf{Findings.}
\begin{itemize}
\item \topEq\ is the most vulnerable crop under both country weightings, in \robTopFfour\ of \robNFfour\ robustness
rankings and in \mcTopShare\% of random weightings. The order below first place is not stable.
\item Water deficit dominates today's scores (\presentWaterShare\%); heat dominates the projected increase. The
flowering-season effective temperature rises by \meanDTeff\,\si{\celsius} and the score rises in \nPairsUp\ of
\nPairs\ pairs.
\item Iron is the nutrient closest to shortfall (\nIronShort\ of \nAdqCountries\ countries under the main
requirement, none under a 15\%-bioavailability requirement).
\item India's ranking is led by wheat, chickpea and common bean; the nine crops supply \shareNewKcal\% of its energy
and \shareNewIron\% of its iron. Hyderabad's rice-based diet coexists with anaemia in \hydAnaemNPW\% of non-pregnant
women in the district; its nutritional risk is more structural than climatic.
\end{itemize}

\noindent\textbf{Four extensions test the foundations.}
\begin{itemize}
\item \textbf{Parameters.} Of \parNRows\ audited parameter rows, the rice and rapeseed heat slopes are unsupported.
Wheat is still first among the seven crops with supported parameters; rice and rapeseed cannot be placed.
\item \textbf{Literature.} \litNRows\ rows from \litNStudies\ publications cover all nine crops; direction agrees where
studies are most comparable and differs where they include CO$_2$ fertilisation or adaptation.
\item \textbf{CO$_2$ and quality.} As a stress test with partial coverage (not a forecast), lower grain concentrations
cut India's zinc supply by \cotwoIndQualZincAbs\% and, with the climate yield loss, by \cotwoIndCombZincAbs\%.
\item \textbf{Diversity.} Losing part of a dominant staple lowers supply but can raise the effective number of food
sources, so diversity must be read with adequacy.
\end{itemize}

\noindent\textbf{What the evidence supports.} The baseline results were regenerated from checksummed raw inputs and are
identical to the previous outputs; \checksPassed\ of \checksN\ consistency checks and \extChecksPassed\ of
\extChecksN\ extension checks pass. This shows that the numbers follow from the stated method, not that the method is
valid. Several parameters are fits or proxies, literature and empirical future yield responses disagree for about half
of the pairs, and validation against yield shocks shows an association only.
\end{abstract}
